\chapter{Tangent space}
\label{ch:tangent}

Equation~\eqref{eq:nts} gives the bumped normal in coordinates relative to the surface. To
light the surface we need it in the same coordinate system as the light and the camera. This
chapter is about the frame that connects the two, and about transforming normals in general.

\section{The tangent frame}

\begin{figure}[ht]
\centering
\begin{tikzpicture}[scale=1.6,>={Stealth}]
  % a curved patch drawn in a simple oblique projection
  \shade[top color=black!3,bottom color=black!12]
    (0,0) .. controls (1,0.35) and (2,0.35) .. (3,0)
    -- (3.8,1.1) .. controls (2.8,1.45) and (1.8,1.45) .. (0.8,1.1) -- cycle;
  \draw (0,0) .. controls (1,0.35) and (2,0.35) .. (3,0)
    -- (3.8,1.1) .. controls (2.8,1.45) and (1.8,1.45) .. (0.8,1.1) -- cycle;
  \coordinate (p) at (1.9,0.68);
  \fill (p) circle (1.2pt);
  \draw[->,very thick,tcolor] (p) -- ++(0.9,0.02) node[right] {$\Tv=\vect P_u/\lVert\vect P_u\rVert$};
  \draw[->,very thick,bcolor] (p) -- ++(-0.38,-0.55) node[below] {$\Bv$ (direction of $+v$)};
  \draw[->,very thick,ncolor] (p) -- ++(0,1.0) node[above] {$\Nv$};
  \draw[->,very thick,npcolor] (p) -- ++(-0.35,0.92) node[left] {$\vect n' = \Tv n_x+\Bv n_y+\Nv n_z$};
\end{tikzpicture}
\caption{The tangent frame at a point of a surface: $\Tv$ along increasing $u$, $\Bv$
along increasing $v$, $\Nv$ outwards. A tangent-space vector $(n_x,n_y,n_z)$ is the
combination $n_x\Tv+n_y\Bv+n_z\Nv$.}
\label{fig:tbn}
\end{figure}

At each point of a textured surface we attach three unit vectors (Figure~\ref{fig:tbn}):
the \emph{tangent} $\Tv$, in the direction of increasing $u$; the \emph{bitangent} $\Bv$
(sometimes, inaccurately, called the binormal), in the direction of increasing $v$; and the
normal $\Nv$. A vector with tangent-space coordinates $\nts=(n_x,n_y,n_z)$ is, in object
coordinates,
\[
  \vect n_{\mathrm{obj}} = n_x\,\Tv + n_y\,\Bv + n_z\,\Nv =
  \underbrace{\begin{pmatrix}\Tv & \Bv & \Nv\end{pmatrix}}_{\text{TBN}}
  \begin{pmatrix}n_x\\ n_y\\ n_z\end{pmatrix}.
\]
This is the famous \emph{TBN matrix}. When its columns are orthonormal, its inverse is its
transpose, so going back---from object space to tangent space, as some techniques do with the
light vector---costs three dot products.

The tangent frame is the reason a single bump or normal map works on any shape: the map
describes the relief \emph{relative to the surface}, and the frame carries it onto whatever
surface it is applied to.

\section{Building an orthonormal frame}

The derivatives $\vect P_u$ and $\vect P_v$ of a parametrization are tangent to the surface but
in general neither unit nor perpendicular (think of the pyramid, Chapter~\ref{ch:texspace}).
Texture Explorer builds an orthonormal frame by the \emph{Gram--Schmidt} process:
\begin{align*}
  \Nv &= \text{the unit outward normal},\\
  \Tv &= \frac{\vect P_u - (\vect P_u\cdot\Nv)\,\Nv}{\lVert\vect P_u - (\vect P_u\cdot\Nv)\,\Nv\rVert},\\
  \Bv &= \sigma\,(\Nv\times\Tv),\qquad \sigma = \operatorname{sign}\bigl((\Nv\times\Tv)\cdot\vect P_v\bigr).
\end{align*}
The sign $\sigma$---stored as the fourth component $w$ of the vertex tangent---records the
\emph{handedness} of the frame. It is needed because a texture can be applied mirrored; with
the conventions of this program, \eqref{eq:orientation} gives $\Tv\times\Bv=-\Nv$ and so
$\sigma=-1$ everywhere. Storing $\Tv$ and $\sigma$ instead of $\Tv$ and $\Bv$ saves two floats
per vertex.

\paragraph{After interpolation.} The rasterizer interpolates $\Nv$ and $\Tv$ linearly across
each triangle. Interpolated unit vectors are shorter than one and no longer perpendicular, so the
pixel shader repeats the Gram--Schmidt step before using them. In production, normal maps are
baked against a precisely specified tangent frame---the \emph{MikkTSpace} convention
\cite{mikkelsen2008}---and the renderer must reproduce exactly the same frame, or the baked
normals come out slightly wrong.

\section{A worked example}
Take the sphere of Chapter~\ref{ch:texspace} at the centre of the texture, $u=v=\frac12$, so
$\theta=\pi/2$ and $\varphi=\pi$. Then
\begin{align*}
  \vect P &= (-1,0,0) = \Nv,\\
  \Tv &= (-\sin\varphi,\,0,\,-\cos\varphi) = (0,0,1),\\
  \Bv &= (\cos\theta\cos\varphi,\,-\sin\theta,\,-\cos\theta\sin\varphi) = (0,-1,0).
\end{align*}
Check the orientation: $\Tv\times\Bv = (0,0,1)\times(0,-1,0) = (1,0,0) = -\Nv$. A tangent-space
normal $\nts = (0.1,\,0,\,0.995)$, leaning towards $+u$, becomes
\[
  \vect n' = 0.1\,\Tv + 0.995\,\Nv = (-0.995,\ 0,\ 0.1)
\]
in object space: it leans towards $+z$, the direction in which $u$ grows at that point. This is
exactly what the Inspector shows when the probe is placed there (with the solid not rotated).

\section{Transforming normals}
\label{sec:normaltransform}

Points and tangent vectors of an object are transformed to the world by its world matrix $M$
(more precisely, by its upper-left $3\times3$ part for vectors). Normals are different. A normal
is defined by being perpendicular to all tangent vectors $\vect t$: $\vect n\cdot\vect t = 0$.
If tangents are transformed as $\vect t' = M\vect t$, the transformed normal must satisfy
$\vect n'\cdot M\vect t = 0$ for every $\vect t$, which holds for
\[
  \vect n' = M^{-\mathsf T}\vect n,\qquad\text{since}\qquad
  (M^{-\mathsf T}\vect n)\cdot(M\vect t) = \vect n^{\mathsf T}M^{-1}M\vect t = \vect n\cdot\vect t = 0 .
\]
For a non-uniform scale, $M^{-\mathsf T}\neq M$ and using $M$ for normals is a classic bug: the
normals of a squashed sphere would lean the wrong way. For a \emph{rotation},
$M^{-\mathsf T}=M$, and for a uniform scale they differ only by a factor that normalization
removes. Texture Explorer's world matrix is a pure rotation about the vertical axis (the
rotation animation), so it transforms points, tangents and normals with the same matrix.

\begin{tryit}
Pin the probe, then start the rotation of the object (space bar). The world-space vectors in
the Inspector change continuously, while the tangent-space normal stays the same: the relief is
a property of the surface, the world-space normal a property of the surface \emph{and} its
placement in the world.
\end{tryit}
